\exercise{Fixation probability}{3}
In a flowerbed there is space for five plants. Plants die at a rate of 3 per plant per time. When this happens, one of the other plants has space to reproduce, so the space freed by the dying plant is filled (immediately) by an offspring of one of the other plants, chosen at random. But one plant has a rare mutation that makes it more resilient such that this plant only dies at rate 2. All of the mutant plant's offspring have the same mutation, whereas neither of the offspring of other plants has it. Compute the probability that mutant plants take over the whole flowerbed.    

\solution 
In this case our system has 6 microstates when we have 0, 1, 2, 3, 4, or 5 mutants. Interestingly there are two absorbing states: If there are 0 mutant plants the mutation is lost and won't come back. If there are 5 mutant plants, then the non-mutant (wildtype) plants have been excluded from the flower bed and the mutation has become fixed in the population. 

We define $x_n$ as the probability that the system is in state $n$ that has $n$ mutant plants. We collect the $x_n$ in the vector 
\eq{
\vec{x} = (x_0,x_1,x_2,x_3,x_4,x_5)^{\rm T}
}
From the question we know that the initial state is 
\eq{
\vec{x}(0)=(0,1,0,0,0,0)^{\rm T}
}
Now there is a hard way and an elegant way to find the answer.

\paragraph{Continuous-time dynamics}
We start by constructing the differential equations. Let's consider first all the processes that reduce the number of mutant plants:
\begin{itemize}
\item If we are in state 1, the mutant plant can die (rate 2) and is then replaced by the wildtype.
\eq{j_{1\to 0} = 2x_1.}
\item If we are in state 2, one of the two mutant plants can die (rate 4), and is then replaced by one of the three wildtype plants (probability 3/4),
\eq{j_{2\to 1} = 4x_2 \frac{3}{4} = 3 x_2.}
\item In state 3 one of the three mutants can die (rate 6) and be replaced by offspring of one of the two wildtype plants (probability 1/2),
\eq{j_{3\to 2} = 6x_3 \frac{1}{2} = 3 x_3.}
\item In state 4, one of the 4 mutants can die (rate 8) and be replaced by offspring from the one wildtype plant (probability 1/4),
\eq{j_{4\to 3} = 8x_4\frac{1}{4} = 2x_4.}
\end{itemize} 
Likewise we can enumerate all the fluxes that increase the number of mutants. 
\begin{itemize}
\item In state 1 one of the four wildtype plants can die (rate 12) and be replaced by offspring of the mutant (probability 1/4)
\eq{j_{1\to 2} = 12x_1\frac{1}{4} = 3x_1.}
\item In state 2 one of the three wildtype plants can die (rate 9) and be replaced by offspring of one of the two mutants (probability 1/2)
\eq{
j_{2\to 3} = 9x_2\frac{1}{2} =4.5x_2
}
\item In state 3 one of the two wildtype plants can die (rate 6) and be replaced by offspring from one of the three mutants (probability 3/4) 
\eq{
j_{3\to 4} = 6x_3\frac{3}{4} =4.5x_3
}
\item In state 4 the one remaining wildtype plant can die (rate 3).
\eq{
j_{4\to 5} = 3x_4
}
\end{itemize}
In summary this gives us the differential equations 
\eqa{
\dot{x}_0 &=& j_{1\to 0}   \\ 
\dot{x}_1 &=& -(j_{1\to 0}+j_{1\to 2}) + j_{2\to 1}    \\  
\dot{x}_2 &=& -(j_{2\to 1}+j_{2\to 3}) + j_{1\to 2} + j_{3 \to 2}  \\  
\dot{x}_3 &=& -(j_{3\to 2}+j_{3\to 4}) + j_{2\to 3} + j_{4 \to 3}   \\  
\dot{x}_4 &=& -(j_{4\to 3}+j_{4\to 5}) + j_{3\to 4}    \\  
\dot{x}_5 &=&  j_{4\to 5}   \\  
}
I don't like having to deal with the fractions, so let's change the unit of time so that everything is twice as fast, then the equations are
\eqa{
\dot{x}_0 &=& 4x_1   \\ 
\dot{x}_1 &=& -(4x_1+6x_1) + 6x_2    \\  
\dot{x}_2 &=& -(6x_2+9x_2) + 6x_1 + 6x_3 \\  
\dot{x}_3 &=& -(6x_3+9x_3) + 9x_2 + 4x_4   \\  
\dot{x}_4 &=& -(4x_4+6x_4) + 9x_3    \\  
\dot{x}_5 &=&  6x_4   \\  
}
which we can also capture by the Jacobian
\eq{
{\bf J} = \avecccccc{
 0 &  4 & 0   & 0   & 0  & 0\\
 0 & -10 & 6  & 0   & 0  & 0\\
 0 &  6 & -15 & 6   & 0  & 0\\
 0 &  0 & 9   & -15 & 4  & 0\\
 0 &  0 &   0 &  9  & -10 & 0 \\ 
 0 &  0 &   0 &  0  & 6  & 0}
}
Now the question is how do we solve for the long term behavior? Clearly the answer will depend on the initial condition and thus structural approaches like the matrix-tree theorem won't help us here. We have to solve the actual dynamics. Since the system is linear the solutions will have the form
\eq{
\vec{x}(t) = \sum_n c_n \vec{v_n} \exp{\lambda_n}  
}
Because our Jacobian is still a Laplacian matrix we know that there can be no positive eigenvalues, however there are two zero eigenvalues. On the one hand we know this from the absorbing states: If there are two absorbing states there must also be two zero eigenvalues. On the other hand we can see it directly in the Jacobian matrix, from the two columns that contain only zeroes.

We can choose the corresponding eigenvectors to be
\eq{
\vec{v_1}=\avec{1\\ 0 \\ 0 \\ 0 \\ 0 \\ 0},\qqq \vec{v_2}=\avec{0\\ 0 \\ 0 \\ 0 \\ 0 \\ 1}
}
Because only two of the eigenvalues are zero, while the others are negative in the limit $t\to \infty$ the general solution will become 
\eq{
\lim_{t\to \infty} \vec{x}(t) = c_1 \vec{v_1} + c_2 \vec{v_2} 
}
If we now consider this solution and the eigenvectors, we can see that $c_1$ is the probability that the mutants go extinct and $c_2$ is the probability that they manage to take over the flowerbed. 

So our task becomes to find the expansion coefficients $c_1$ and $c_2$. The easiest way to determine these coefficents is the formula 
\eq{
c_n = \vec{w_n}^{\dag}\vec{x}(0)
}
where $\vec{w_n}^\dag$ is the $n$-th left eigenvector, which has been normalized such that 
\eq{
\vec{w_n}^{\dag}\vec{v_n}=1
}

We thus have to solve 
\eqa{
 0\vec{w}^\dag &=& \vec{w}^\dag {\bf J} \\
   &=& (w_0,w_1,w_2,w_3,w_4,w_5) \avecccccc{
 0 &  4 & 0   & 0   & 0  & 0\\
 0 & -10 & 6  & 0   & 0  & 0\\
 0 &  6 & -15 & 6   & 0  & 0\\
 0 &  0 & 9   & -15 & 4  & 0\\
 0 &  0 &   0 &  9  & -10 & 0 \\ 
 0 &  0 &   0 &  0  & 6  & 0}
} 
From the product of the vector with the non-zero columns we get the conditions  
\eqa{
0&=& 4w_0 -10w_1 + 6w_2 \\
0&=& 6w_1 -15w_2 + 9w_3 \\
0&=& 6w_2 -15w_3 + 9w_4 \\
0&=& 4w_3 -10w_4  +6w_5 
}
So we have four conditions for 6 unknowns, which means that we can chose two of the unknowns as we want. So how should we chose them? Recall from Chap.~20 that our method for computing the expansion coefficients requires that the left eigenvectors are orthogonal to all right eigenvectors that belong to different eigenvalues. In particular this means that when we are constructing $\vec{w_1}^\dag$ then we want to make sure that it is orthogonal to $\vec{v_2}$. Because $\vec{v_2}$ only has one non-zero element and it is in the six row any vector that obeys the orthogonality condition $\vec{w}^\dag\vec{v_2}=0$ must have $w_5=0$, so this is the first choice that we make. Moreover we want our vector to be normalized such that $\vec{w_1}^\dag\vec{v_1}=1$ which we can ensure by choosing $w_0=1$. 

Taking these constraints into account the equation system reads 
\eqa{
0&=& 4 -10w_1 + 6w_2 \\
0&=& 6w_1 -15w_2 + 9w_3 \\
0&=& 6w_2 -15w_3 + 9w_4 \\
0&=& 4w_3 -10w_4   
}
To simplify it is useful to divide the first and last condition by 2 and the other two conditions by 3. 
\eqa{
0&=& 2 -5w_1 + 3w_2 \\
0&=& 2w_1 -5w_2 + 3w_3 \\
0&=& 2w_2 -5w_3 + 3w_4 \\
0&=& 2w_3 -5w_4   
}
We can now solve the forth condition for 
\eq{
w_4= \frac{2}{5} w_3.
}
Substituting back leaves us with 
\eqa{
0&=& 2 -5w_1 + 3w_2 \\
0&=& 2w_1 -5w_2 + 3w_3 \\
0&=& 2w_2 -5w_3 + 6w_3/5 
}
Multiplying the third condition by 5 yields 
\eq{
0=10w_2 -25w_3 +6w_3
}
and hence 
\eq{
w_2 = \frac{19}{10} w_3
}
So now we only have two conditions left to satisfy 
\eqa{
0&=& 2 -5w_1 + 57w_3/10 \\
0&=& 2w_1 -95 w_3/10 + 3w_3 
}
We can write the second condition as  
\eq{
0= 4w_1 - 19 w_3 + 6 w_2
}
which we can solve for 
\eq{
w_1 = \frac{13}{4} w_3
}
The only remaining condition is now 
\eq{
0=2 - \frac{65}{4} w_3 + \frac{57}{10} w_3
}
which we multiply by 20 to find 
\eq{
0=40 - 325 w_3 + 114 w_3 
}
which we solve for
\eq{
w_3=\frac{40}{211}
}
The element that we really need is $w_1$ so we use the relationship from above to find
\eq{
w_1 = \frac{13}{4} w_3 = \frac{13}{4} \frac{40}{211} = \frac{130}{211} \approx 61.6\%
}

Let's summarize: We know that in the long run, the state of the flowerbed is
\eq{
\lim_{t\to \infty} \vec{x}(t) = c_1 \vec{v_1} + c_2 \vec{v_2} 
}
where $v_1$ is the state in which the mutant becomes extinct. We can compute $c_1$ as 
\eq{
c_1 = \vec{w^\dag} \vec{x(0)} = (w_0,w_1,w_2,w_3,w_4,w_5)\avec{0\\ 1\\ 0 \\ 0 \\ 0 \\ 0} = w_1 \approx 61.6\%
}
Hence the probability that the muntant goes extinct is approximately 61.6\% and the probability that it takes over (i.e.~the mutation becomes fixed in the population) is complementary probability of approximately 38.4\%. 

\paragraph{Discrete-time dynamics}  When we computed the Jacobian in the solution above you may have realized that in all non-absorbing states the following is true: The next state has one more mutant with probability 3/5 and one less mutant with probability 2/5. So we can also interpret the system as a discrete time random walk with the propagator 
\eq{
{\bf P} = \avecccccc{ 
   1 & 2/5 & 0   & 0   & 0    &  0 \\
   0 &  0  & 2/5 & 0   & 0    &  0 \\
   0 & 3/5 &  0  & 2/5 & 0    &  0 \\
   0 &  0  & 3/5 & 0   & 2/5  &  0 \\
   0 &  0  &  0  & 3/5 &  0   &  0 \\
   0 &  0  &  0  &  0  & 3/5  &  1  }
}
where we put a 1 on the diagonal in the first and last row to indicate that when we land in the first or last state then the 
system is stuck there forever. 

Because we now have a discrete-time system the long term behaviour $\vec{x}$ is captured by the condition 
\eq{
\vec{x} = {\bf P} \vec{x}
}
Thus we are now looking at eigenvectors of $\bf P$ that have the eigenvalue 1. From our understanding of the system we can guess that these must be the eigenvectors 
\eq{
 \vec{v_1} = \avec{1 \\ 0 \\ 0 \\ 0 \\ 0  \\ 0} \qqq \vec{v_2} = \avec{0 \\ 0 \\ 0 \\ 0 \\ 0  \\ 1}
}
That is either we end up with zero mutants ($\vec{v_1}$) or 5 mutants ($\vec{v_2}$). Again the long term behaviour is composed from these two scenarios such that  
\eq{
\vec{\vec{x}}(\infty) = c_1 \vec{v_1} + c_2 \vec{v_2} 
}
As in the continuous case we can now find the expansion coefficients as 
\eq{
c_n = \vec{w_n^\dag} \vec{x}(0) 
}
where $\vec{w_n^\dag}$ are the left eigenvectors. They obey the condition 
\eq{
\vec{w^\dag}= \vec{w^\dag} {\bf P} 
}
and thus 
\eqa{
w_0 &=& w_0 \\
w_1 &=& \frac{2}{5} w_0 + \frac{3}{5} w_2 \\
w_2 &=& \frac{2}{5} w_1 + \frac{3}{5} w_3 \\
w_3 &=& \frac{2}{5} w_2 + \frac{3}{5} w_4 \\
w_4 &=& \frac{2}{5} w_3 + \frac{3}{5} w_5 \\
w_5 &=& w_5
}
Again we get two trivial conditions and four nontrivial ones. So we can now eliminate the two trivial conditions. Let's say we want the eigenvector $\vec{w_2^\dag}$ this time. So this eigenvector must obey 
\eq{
   \vec{w_2^\dag}\vec{v_1}=0 \qqq \vec{w_2^\dag}\vec{v_2}=1 
}
which tells us $w_0=0$ and $w_5=1$. The quickest way to go is to ignore the normalization condition $w_5=1$ for now. Instead we will just try $w_1=1$ and if we end up with a $w_5$ that is too high or too low then we will just rescale the vector accordingly. Substituting $w_0=0$, $w_1=1$ into the condition for $w_1$ yields 
\eq{
1 = 0 + \frac{3}{5} w_2 
}
which means 
\eq{
w_2=\frac{5}{3}
}
Substituting into the condition for $w_2$ gives us 
\eq{
\frac{5}{3} = \frac{2}{5} + \frac{3}{5}w_3 
}
and hence 
\eq{
w_3 = \frac{25}{9} - \frac{2}{3} = \frac{19}{9}
}
Similar the condition for $w_3$ gives us 
\eq{
\frac{19}{9} = \frac{2}{5}\frac{5}{3} + \frac{3}{5}w_4  
}
and hence 
\eq{
w_4 = \frac{95}{27} - \frac{10}{9} = \frac{65}{27}.
}
Finally the conditions for $w_4$ tells us
\eq{
\frac{65}{27} = \frac{2}{5}\frac{19}{9} + \frac{3}{5} w_5
}
and therefore 
\eq{
w_5 = \frac{325}{81} - \frac{38}{27} = \frac{211}{81}
}
So instead of the desired value of $w_5=1$ we got $211/81$. Because everything is nicely linear our assumed value for $w_1$ was too high by the same factor, which tells us that the true value is 
\eq{
w_1= \frac{81}{211} \approx 38.4\%
}
So this our result. As above we find a fixation probability of 38.4\% for the mutation. 


